\documentclass[11pt,a4paper]{article}

\usepackage[utf8]{inputenc}
\usepackage[margin=0.85in]{geometry}
\usepackage{amsmath,amssymb,amsfonts}
\usepackage{graphicx}
\usepackage{subcaption}
\usepackage{float}
\usepackage{hyperref}
\usepackage{microtype}

\title{\textbf{\Large Part 4: Frequency-Domain Filtering and Image Restoration}}
\author{}
\date{}

\begin{document}

\maketitle

% =========================================================
\section{Motion Blur Estimation and Transfer Function}
% =========================================================

\subsection{Estimating the Blur Vector}

The image was captured while rotating the camera, producing predominantly horizontal motion blur. We model the motion as

\[
x_0(t)=at,\qquad y_0(t)=bt,\qquad t\in[0,1].
\]

By examining the length and direction of the visible streaks, we estimated the blur to be approximately 30 pixels horizontally, with negligible vertical motion. Thus,

\[
\boxed{(a,b)=(30,0)\text{ pixels}.}
\]

\subsection{Transfer Function}

The blurred image is

\[
g(x,y)=\int_0^1 f(x-at,y-bt)\,dt.
\]

Taking the Fourier transform gives

\[
G(u,v)=H(u,v)F(u,v),
\]

where

\[
H(u,v)
=
\int_0^1 e^{-i2\pi(ua+vb)t}\,dt.
\]

Therefore,

\[
\boxed{
H(u,v)
=
e^{-i\pi(ua+vb)}
\operatorname{sinc}(ua+vb)
}
\]

where

\[
\operatorname{sinc}(x)=\frac{\sin(\pi x)}{\pi x}.
\]

For $(a,b)=(30,0)$,

\[
\boxed{
H(u,v)=e^{-i30\pi u}\operatorname{sinc}(30u)
}
\]

and hence

\[
\boxed{|H(u,v)|=|\operatorname{sinc}(30u)|.}
\]

Since the blur is horizontal, the zeros of $H$ form approximately vertical bands in the frequency domain.

\begin{figure}[H]
    \centering
    \begin{subfigure}[b]{0.45\textwidth}
        \includegraphics[width=\textwidth]{01_input.png}
        \caption{Input blurred image}
    \end{subfigure}
    \hfill
    \begin{subfigure}[b]{0.45\textwidth}
        \includegraphics[width=\textwidth]{02_H_magnitude.png}
        \caption{$|H(u,v)|$}
    \end{subfigure}
    \caption{Input image and magnitude of the motion-blur transfer function.}
\end{figure}

The Fourier spectrum shows the expected directional structure caused by the horizontal motion blur.

% =========================================================
\section{Inverse Filtering}
% =========================================================

Inverse filtering estimates the original spectrum as

\begin{figure}[H]
    \centering
    \includegraphics[width=0.35\linewidth]{03_inverse_filter.png}
    \caption{Result of inverse filtering.}
\end{figure}
\[
\boxed{
\hat F(u,v)=\frac{G(u,v)}{H(u,v)}.
}
\]

Since $H$ has zeros, direct division is unstable. In the implementation, values satisfying

\[
|H(u,v)|<10^{-3}
\]

were replaced by $10^{-3}$.



The result contains severe streaking artifacts. This occurs because frequencies for which $|H|$ is small are strongly amplified by $1/H$, including noise and modelling errors. Since the problematic frequencies occur along specific bands rather than uniformly in the frequency plane, the resulting artifacts are directional rather than isotropic.

% =========================================================
\section{Wiener Deconvolution}
% =========================================================

To reduce the instability of inverse filtering, we used the Wiener filter

\[
\boxed{
\hat F(u,v)
=
\frac{H^*(u,v)}
{|H(u,v)|^2+K}G(u,v).
}
\]

We tested a range of $K$ values. Small values produced strong streaking, while large values excessively smoothed the image. Based on visual comparison, we selected

\[
\boxed{K_{\mathrm{best}}=0.005}.
\]

The required comparison is therefore

\[
0.1K=0.0005,\qquad
K=0.005,\qquad
10K=0.05.
\]

\begin{figure}[H]
    \centering
    \begin{subfigure}[b]{0.32\textwidth}
        \includegraphics[width=\textwidth]{04_wiener_0.1K.png}
        \caption{$K=0.0005$}
    \end{subfigure}
    \hfill
    \begin{subfigure}[b]{0.32\textwidth}
        \includegraphics[width=\textwidth]{05_wiener_best.png}
        \caption{$K=0.005$}
    \end{subfigure}
    \hfill
    \begin{subfigure}[b]{0.32\textwidth}
        \includegraphics[width=\textwidth]{06_wiener_10K.png}
        \caption{$K=0.05$}
    \end{subfigure}
    \caption{Wiener deconvolution for $0.1K$, $K$, and $10K$.}
\end{figure}

For $K=0.0005$, the restoration is extremely aggressive, but streaking noise completely dominates the image. At $K=0.005$, the image retains strong sharpness, but there are still distinct line artifacts present. These lines are ringing artifacts caused by the near-zero values of the transfer function $H(u,v)$ in the frequency domain. There is an inherent trade-off here: if we try to fully remove these lines by increasing the regularization further (e.g., to $K=0.05$), the artifacts successfully vanish, but the high-frequency details are also heavily suppressed, causing the overall image to become unacceptably blurry. Thus, $K=0.005$ gives the best qualitative compromise by maximizing the deblurring sharpness at the cost of retaining some of these artifact lines.

% =========================================================
\section{Regularized Deconvolution}
% =========================================================

\begin{figure}[H]
    \centering
    \begin{subfigure}[b]{0.32\textwidth}
        \includegraphics[width=\textwidth]{07_reg_0.1lam.png}
        \caption{$0.1\lambda=0.005$}
    \end{subfigure}
    \hfill
    \begin{subfigure}[b]{0.32\textwidth}
        \includegraphics[width=\textwidth]{08_reg_best.png}
        \caption{$\lambda=0.05$}
    \end{subfigure}
    \hfill
    \begin{subfigure}[b]{0.32\textwidth}
        \includegraphics[width=\textwidth]{09_reg_10lam.png}
        \caption{$10\lambda=0.5$}
    \end{subfigure}
    \caption{Regularized deconvolution for three values of $\lambda$.}
\end{figure}

We use the squared gradient norm specified in the assignment:

\[
R(\hat f)
=
\sum_{x,y}
\left[
(\hat f(x+1,y)-\hat f(x,y))^2+
(\hat f(x,y+1)-\hat f(x,y))^2
\right].
\]

The objective is

\[
E(\hat f)=\|h*\hat f-g\|^2+\lambda R(\hat f).
\]

For the finite differences,

\[
|D_x(u,v)|^2=4\sin^2(\pi u),
\qquad
|D_y(u,v)|^2=4\sin^2(\pi v).
\]

Thus,

\[
P(u,v)
=
4\sin^2(\pi u)+4\sin^2(\pi v).
\]

In the Fourier domain,

\[
E(\hat F)
=
\sum_{u,v}
\left[
|H\hat F-G|^2+
\lambda P|\hat F|^2
\right].
\]

Setting the derivative with respect to $\hat F^*$ to zero gives

\[
\boxed{
\hat F(u,v)
=
\frac{H^*(u,v)G(u,v)}
{|H(u,v)|^2+\lambda P(u,v)}
}
\]

or equivalently,

\[
\boxed{
\hat F(u,v)
=
\frac{H^*(u,v)G(u,v)}
{|H(u,v)|^2+
\lambda\left[
4\sin^2(\pi u)+4\sin^2(\pi v)
\right]}.
}
\]

We tested a wide range of $\lambda$ values and found

\[
\boxed{\lambda_{\mathrm{best}}=0.05}
\]

to give the best qualitative result.



With $\lambda=0.005$, more streaking remains due to weak regularization. At $\lambda=0.05$, the artifacts are substantially reduced while image structure is retained. For $\lambda=0.5$, stronger smoothing removes additional high-frequency detail. Hence, $\lambda=0.05$ was selected.





\end{document}